% $Id: template.tex 11 2007-04-03 22:25:53Z jpeltier $

\documentclass{vgtc}                          % final (conference style)
% \documentclass[review]{vgtc}                 % review
%\documentclass[widereview]{vgtc}             % wide-spaced review
%\documentclass[preprint]{vgtc}               % preprint
%\documentclass[electronic]{vgtc}             % electronic version

%% Uncomment one of the lines above depending on where your paper is
%% in the conference process. ``review'' and ``widereview'' are for review
%% submission, ``preprint'' is for pre-publication, and the final version
%% doesn't use a specific qualifier. Further, ``electronic'' includes
%% hyperreferences for more convenient online viewing.

%% Please use one of the ``review'' options in combination with the
%% assigned online id (see below) ONLY if your paper uses a double blind
%% review process. Some conferences, like IEEE Vis and InfoVis, have NOT
%% in the past.

%% Figures should be in CMYK or Grey scale format, otherwise, colour 
%% shifting may occur during the printing process.

%% it is recomended to use ``\cref{sec:bla}'' instead of ``Fig.~\ref{sec:bla}''
\graphicspath{{figures/}{pictures/}{images/}{./}} % where to search for the images

\usepackage{times}                     % we use Times as the main font
\renewcommand*\ttdefault{txtt}         % a nicer typewriter font

%% Only used in the template examples. You can remove these lines.
\usepackage{tabu}                      % only used for the table example
\usepackage{booktabs}                  % only used for the table example
\usepackage{lipsum}                    % used to generate placeholder text
\usepackage{mwe}                       % used to generate placeholder figures

%% We encourage the use of mathptmx for consistent usage of times font
%% throughout the proceedings. However, if you encounter conflicts
%% with other math-related packages, you may want to disable it.
\usepackage{mathptmx}                  % use matching math font
\usepackage{amsmath}
\usepackage{amssymb}
\usepackage{array}
\usepackage{multirow}
\usepackage{graphicx}
\usepackage{xcolor}
\usepackage{url}
\usepackage{enumitem}

\newcommand{\R}{\mathbb{R}}
\newcommand{\Domain}{\mathcal{M}}
\newcommand{\Range}{\mathcal{R}}
\newcommand{\Sheets}{\mathcal{S}}
\newcommand{\Sheet}{S}
\newcommand{\Score}{C}
\newcommand{\placeholder}[2]{%
  \begin{center}
  \fbox{%
    \begin{minipage}[c][#1][c]{0.92\linewidth}
      \centering
      #2
    \end{minipage}}
  \end{center}
}
\newcommand{\wideplaceholder}[2]{%
  \begin{center}
  \fbox{%
    \begin{minipage}[c][#1][c]{0.96\textwidth}
      \centering
      #2
    \end{minipage}}
  \end{center}
}
\newcommand{\myparagraph}[1]{\vspace{1mm} \noindent \textbf{#1}}
\newcommand{\ie}{i.e.,\xspace}
\newcommand{\etal}{et al.\xspace}
\newcommand{\Mohit}[1]{\textcolor{blue}{#1}\xspace}
\newcommand{\Talha}[1]{\textcolor{magenta}{Talha: #1}\xspace}
\newcommand{\Nanna}[1]{\textcolor{purple}{#1}\xspace}
\newcommand{\Ingrid}[1]{\textcolor{red}{#1}\xspace}
\newcommand{\Petar}[1]{\textcolor{cyan}{#1}\xspace}

%% If you are submitting a paper to a conference for review with a double
%% blind reviewing process, please replace the value ``0'' below with your
%% OnlineID. Otherwise, you may safely leave it at ``0''.
\onlineid{0}

%% declare the category of your paper, only shown in review mode
\vgtccategory{Research}

%% allow for this line if you want the electronic option to work properly
\vgtcinsertpkg

%% In preprint mode you may define your own headline. If not, the default IEEE copyright message will appear in preprint mode.
%\preprinttext{To appear in an IEEE VGTC sponsored conference.}

%% This adds a link to the version of the paper on IEEEXplore
%% Uncomment this line when you produce a preprint version of the article 
%% after the article receives a DOI for the paper from IEEE
%\ieeedoi{xx.xxxx/TVCG.201x.xxxxxxx}


%% Paper title.

\title{Topological Tracking of Bivariate Features Using Reeb Spaces}

%% This is how authors are specified in the conference style

%% Author and Affiliation (single author).
%%\author{Roy G. Biv\thanks{e-mail: roy.g.biv@aol.com}}
%%\affiliation{\scriptsize Allied Widgets Research}

%% Author and Affiliation (multiple authors with single affiliations).
%%\author{Roy G. Biv\thanks{e-mail: roy.g.biv@aol.com} %
%%\and Ed Grimley\thanks{e-mail:ed.grimley@aol.com} %
%%\and Martha Stewart\thanks{e-mail:martha.stewart@marthastewart.com}}
%%\affiliation{\scriptsize Martha Stewart Enterprises \\ Microsoft Research}

%% Author and Affiliation (multiple authors with multiple affiliations)
\author{Mohit Sharma\thanks{e-mail: mohit.sharma@liu.se}\\ %
        \scriptsize Link\"oping university %
\and Petar Hristov\thanks{e-mail: petar.georgiev.hristov@liu.se}\\ %
        \scriptsize Link\"oping university %
\and Nanna Holmgaard List \thanks{e-mail: nalist@kth.se}\\ %
        \scriptsize KTH Royal Institute of Technology %
\and Talha Bin Masood\thanks{e-mail: talha.bin.masood@liu.se}\\ %
        \scriptsize Link\"oping university %
\and Ingrid Hotz\thanks{e-mail: ingrid.hotz@liu.se}\\ %
        \scriptsize Link\"oping university %
}

%% A teaser figure can be included as follows
\teaser{
  \centering
  \includegraphics[width=\linewidth]{CypressView}
  \caption{In the Clouds: Vancouver from Cypress Mountain.}
  \label{fig:teaser}
}

%% Abstract section.
\abstract{
Time-varying bivariate fields arise in many scientific applications, yet their feature-level temporal analysis remains challenging because bivariate features are not captured by scalar-field extrema, level sets, or contour-trees. Reeb spaces provide a topological representation of bivariate fields, but using them for temporal analysis is difficult: computed sheets can be noisy, their range-space projections may overlap, and meaningful correspondences across timesteps are nontrivial to define. Existing continuous-scatterplot-based approaches provide useful timestep-level summaries, including descriptors of domain-segmentation-driven peeled CSP layers, but they do not directly provide a topological feature graph for tracking individual bivariate features. We propose a Reeb space-based framework for tracking bivariate features over time. At each timestep, we compute the Reeb space, extract prominent sheets by area, and estimate correspondences between sheets in adjacent timesteps using complementary similarity measures in the domain and range. We summarize the resulting temporal correspondences using a Sankey-based sheet evolution summary linked with sheet, fiber-surface, and diagnostic views. We evaluate the method on two synthetic torus datasets and three time-varying molecular electronic-structure datasets. The results show that Reeb space sheet tracking reveals persistent feature families, abrupt reconfiguration intervals, and disagreements between domain- and range-based similarities that are difficult to identify from temporal descriptors of peeled CSP layers. Overall, the framework demonstrates that Reeb space sheets can serve as trackable topological features for time-varying bivariate data.



} % end of abstract

%% Keywords that describe your work. Will show as 'Index Terms' in journal
%% please capitalize first letter and insert punctuation after last keyword.
\keywords{Bivariate field, Reeb space, fiber surface, time-varying data, molecular visualization.}

%% Copyright space is enabled by default as required by guidelines.
%% It is disabled by the 'review' option or via the following command:
% \nocopyrightspace

%%%%%%%%%%%%%%%%%%%%%%%%%%%%%%%%%%%%%%%%%%%%%%%%%%%%%%%%%%%%%%%%
%%%%%%%%%%%%%%%%%%%%%% START OF THE PAPER %%%%%%%%%%%%%%%%%%%%%%
%%%%%%%%%%%%%%%%%%%%%%%%%%%%%%%%%%%%%%%%%%%%%%%%%%%%%%%%%%%%%%%%%

\begin{document}

%% The ``\maketitle'' command must be the first command after the
%% ``\begin{document}'' command. It prepares and prints the title block.

%% the only exception to this rule is the \firstsection command
\firstsection{Introduction}

\maketitle
Many scientific simulations produce multiple scalar fields over a common spatial domain. In molecular electronic-structure analysis, for example, an electronic transition can be represented using hole and particle natural transition orbital (NTO) fields~\cite{Martin2003NTO}. These fields indicate where charge is depleted and accumulated during the transition. As the molecular geometry evolves, the analysis becomes a time-varying bivariate problem: the goal is not only to study a pair of fields, but also to understand how their joint features evolve over time.

A common strategy for analyzing bivariate fields is to inspect selected isosurfaces or fiber surfaces. Fiber surfaces generalize isosurfaces to bivariate data and provide detailed local views of structures associated with chosen range values. They are useful for examining specific hypotheses, but they do not by themselves provide a compact representation of temporal feature evolution. Another strategy is to summarize each timestep using a continuous scatterplot (CSP). CSPs generalize discrete scatterplots to continuous spatial fields and summarize the distribution of bivariate values. Recent work has used CSP descriptors to analyze the temporal evolution of molecular electronic-structure data~\cite{Sharma2025}. In that approach, bivariate features are obtained using domain-segmentation-driven CSP peeling~\cite{Sharma2021}. The resulting descriptors capture the temporal behavior of the peeled CSP layers. However, they do not directly provide a topological feature evolution that identifies how individual bivariate features persist, split, disappear, merge, or reconfigure across timesteps.

We address this limitation using Reeb spaces. For a bivariate field $F=(f,g):\Domain\rightarrow\R^2$, the Reeb space is defined by identifying domain points that belong to the same connected component of a fiber $F^{-1}(a,b)$. Its two-dimensional cells, called sheets, capture coherent bivariate structures by combining their extent in the range with connectedness information in the domain. Thus, sheets provide a topological description of how bivariate features are organized.

This paper proposes a practical pipeline for tracking prominent Reeb space sheets in time-varying bivariate data. At each timestep, we compute the Reeb space, extract prominent sheets by area, and generate sheet and fiber-surface renderings for detailed inspection. We then compute correspondences between sheets in adjacent timesteps and summarize the resulting temporal graph using a Sankey-based sheet evolution summary. The tracking uses complementary similarity measures in the range and domain. Range-space similarity measures whether two sheets occupy similar regions in the bivariate range. Domain-space similarity measures whether they are supported by similar regions of the input mesh. These notions capture different aspects of sheet correspondence. Their agreement suggests stable feature evolution, while their disagreement can indicate reconfiguration events.

We evaluate the approach on synthetic torus datasets and time-varying molecular electronic-structure datasets. The synthetic examples provide controlled cases for studying sheet persistence, appearance, disappearance, merging, and splitting. The molecular datasets demonstrate how Reeb space tracking can complement CSP-based temporal analysis by identifying persistent feature families and important transition intervals. Together, these results show that Reeb space sheets can serve as trackable topological features for time-varying bivariate fields. The main contributions of this paper are:
\begin{itemize} 
\item A feature-level temporal representation for time-varying bivariate fields based on correspondences between Reeb space sheets. 
\item Complementary sheet similarity measures that combine range-space overlap with domain-space support overlap. 
\item A visual analysis workflow that links a Sankey-based temporal summary with sheet renderings, fiber-surface renderings, centroid coloring, and diagnostic summaries. 
\item A quantitative and qualitative evaluation on synthetic and molecular datasets, including metric-agreement analysis and event-interval detection. \end{itemize}





\begin{figure*}[t]
    \centering
    \includegraphics[width=0.9\textwidth]{./figures/timer-varying.pdf}
    \caption{Introduction to the main topological and geometric structures used in the paper.
    We assume that our domain \emph{a)} is a tetrahedral mesh and we render a fiber and a fiber surface.
    In the range \emph{b)} we show the arrangement of the images of the edges of the tetrahedra. 
    Combining the faces of this arrangement yields the three sheets of the Reeb space \emph{c)}.
    Finally, the sheets of the Reeb space correspond to three volumes which segment the domain.
    We call those volumes features.
    }
    \label{fig:pipeline}
\end{figure*}

\section{Related work}
\Talha{First draft largely taken from our previous papers. TODO: Update and shorten; Add a few more up-to-date references, and fix issues with bibtex entries.}

\myparagraph{Reeb space.} The Reeb space is a higher-dimensional generalization of the Reeb graph that captures the topology of multiple scalar functions defined on a common domain \cite{Edelsbrunner2008}. Whereas a Reeb graph summarizes the connectivity of level sets of a single scalar field, the Reeb space describes the connectivity of \emph{fibers} in multivariate data. In the bivariate case, it consists of a collection of two-dimensional sheets connected along Jacobi edges \cite{Edelsbrunner2004b}. This structure provides a compact representation of the topological organization of a multifield and forms the basis of many approaches to multivariate topological analysis. Several algorithms have been proposed for computing Reeb spaces in bivariate settings \cite{Tierny2017,chattopadhyayAlgorithmFastCorrect2024,Hristov2026SingularAA}, as well as for more general domains and higher-dimensional ranges \cite{Edelsbrunner2008,hristov2025}. One of the key challenges in these methods is achieving robustness in the presence of numerical inaccuracies \cite{Hristov2025b}. To improve scalability, approximate representations such as the Joint Contour Net \cite{Carr2014} and Mapper \cite{Dey2016,Brown2021} have also been developed, providing some practical alternatives for the analysis of complex multivariate datasets. In addition to serving as topological descriptors, Reeb spaces and their approximations have been used to compare multifields by defining distances between their topological structures \cite{Ramamurthi2022a,Agarwal2021,Ramamurthi2024}.

\myparagraph{Approaches to feature tracking.} Topology-based feature tracking is a central task in topological data analysis and visualization of scalar and vector fields. Broadly, tracking approaches follow two paradigms: either features are extracted independently at each time step and subsequently matched, or the temporal dimension is incorporated directly into the data domain, transforming tracking into a spatio-temporal feature extraction problem \cite{Post2003,Weber2011}. Topological features are commonly represented either by critical points, such as maxima and minima, or by spatial regions defined through sub-, super-, or level sets. Intermediate feature definitions also exist, including groups of critical points represented by merge-tree structures, such as crown features, which have been employed in applications ranging from cyclone tracking to medical imaging \cite{Engelke2020,Nilsson2022Cyclone,Rasheed2022}. Other application-specific descriptors include vortex features derived from scalar quantities in flow fields and extremal structures based on subsets of the Morse-Smale complex.

A large body of work focuses on tracking critical point-based features. Early approaches tracked vortices and other flow phenomena through critical points in derived scalar fields \cite{Garth2004b,Weinkauf2011}. Reininghaus et al.~\cite{Reininghaus2011} introduced Combinatorial Feature Flow Fields, a history-based framework that has since been extended to track hierarchical critical point structures and applied in several domains \cite{Engelke2020,Nilsson2022Cyclone}. More recently, correspondence has been established through topological similarity measures, including persistence-diagram matching, merge-tree mappings, and tree-edit distances \cite{Cohen-Steiner2010,Sridharamurthy2020,Wetzels2022a,Wetzels2023}. Alternative approaches combine topological and geometric information, exploit temporal interpolation, or use optical-flow-based correspondences \cite{Yan2022b,Valsangkar2019}. Critical points have also been tracked indirectly through their associated ascending and descending manifolds in the Morse-Smale complex \cite{Nilsson2023}.

For spatial features, correspondence is typically determined through similarity measures such as spatial distance, attribute similarity, and, most commonly, spatial overlap. Foundational work by Samtaney et al.~\cite{Samtaney1994} and Silver and Wang~\cite{silver1998tracking} established overlap-based tracking of threshold-defined regions, augmented by geometric attributes such as volume and center of mass. Spatial overlap remains widely used, including in methods based on merge trees, where hierarchical feature relationships can be exploited for tracking and visualization through nested tracking graphs \cite{Widanagamaachchi2012,Saikia2017,Lukasczyk2017nested}. Additional processing steps, such as semantic filtering or global optimization, are often employed to remove implausible correspondences and construct coherent feature tracks \cite{Nilsson2022Cyclone,Schnorr2020}. Hierarchical tracking methods leverage the nested structure of level sets and merge trees to capture multi-scale behavior \cite{Lukasczyk2019dynamic}, while layout and visualization techniques have improved the readability of the resulting tracking graphs \cite{Kopp2019,Dobler2024}. In contrast to traditional applications that focus on sparse features occupying a small fraction of the domain, Schnorr et al.~\cite{Schnorr2020} considered space-filling features in combustion simulations, where the entire domain is partitioned into regions. Although overlap-based correspondence remains effective in this setting, it produces highly connected tracking graphs that require computationally expensive optimization procedures to extract meaningful feature trajectories.


\section{Background}
This section introduces the concepts used throughout the paper. We begin with bivariate fields and their domain and range spaces, then review fibers, fiber surfaces, Reeb spaces, and Reeb space sheets. These concepts provide the basis for formulating sheet tracking as a temporal correspondence problem between prominent Reeb space sheets and their corresponding features extracted at consecutive timesteps.

\myparagraph{Bivariate field.} Let $\Domain$ be a triangulation of a regular-grid volume or more generally a 3-manifold.
A \textit{scalar field} over $\Domain$ is a piecewise-linear map $f:\Domain \to \R$, typically defined on the vertices and interpolated along the higher dimensional simplices using barycentric interpolation.
A \textit{bivariate field} is a piecewise-linear map $F=(f,g):\Domain \to \R^2,$
where $f$ and $g$ are scalar fields defined on the same domain.
We refer to $\Domain$ as the \textit{domain space}, to $\R^2$ as the \textit{range space} and to $F(\Domain)$ as the \textit{image} of $\Domain$ under $F$.
A bivariate field is \textit{generic} when the image of no three vertices of $\Domain$ are collinear in the plane and the images of no three edges are concurrent in their interiors.
Generic maps ensure that preimages are well-behaved.

\myparagraph{Fiber.} The preimage of a point $y=(a,b)\in\R^2$ defines a \emph{fiber}
$F^{-1}(y)=\{x \in \Domain \mid f(x)=a,\ g(x)=b\}$.
The fibers of generic bivariate maps defined over three-dimensional domains are typically a collection of one-dimensional curves.
Fiber can only change their connectivity when they intersect singular edges and vertices~\cite{hristovHypersweepsConvectiveClouds2022}.
The preimage $F^{-1}(P)$ of a polygon $P \subset \R^2$ in the range is a collection of fibers that sweep out a \emph{fiber surface}.
Fiber surfaces are useful for interpreting how selected range-space structures are embedded in the spatial domain.


\myparagraph{Reeb space.} The Reeb space generalizes the Reeb graph from scalar functions to multivariate maps~\cite{Edelsbrunner2008}.
It captures how the connected components of fibers change for all points over the range of the map.
The Reeb space is the quotient space $\mathcal{R}_F = \Domain/\sim$, where $\sim$ is the equivalence relation on $\Domain$ such that $x\sim x' \text{ if } F(x) = F(x') \text{ and }
x, x' \text{ lie in the same connected component of } F^{-1}(F(x))$.
Practical algorithms for computing the Reeb space include Jacobi fiber surfaces~\cite{Tierny2017b} and the singular arrange and traverse~\cite{Hristov2026SingularAA}.

\myparagraph{Reeb Space Sheets.} 
For a bivariate map over a three-dimensional domain, the Reeb space can be represented by a two-dimensional cell complex. 
The 2-cells of the Reeb space are referred to as \textit{sheets}.
The sheets of the Reeb space are glued along singular vertices and edges where the connectivity of fibers changes, called the Jacobi structure~\cite{Chattopadhyay2014}.
Preimages of sheets in the domain correspond to features captured by the Reeb space, analogous to how to preimages of edges of the Reeb graph and contour tree correspond to features in scalar fields.

\myparagraph{Sheet and Feature Metrics.}
Each sheets of the Reeb space can be described geometrically by its projection onto the range space as a polygon.
In the arrange and traverse framework~\cite{hristov2025, Hristov2026SingularAA}, each such polygon is the union of faces of the geometric arrangements of the images of singular edges.
The area of each sheet is then the sum of areas of its associated arrangement faces.
The volume of a feature associated with a sheet can be computed exactly using the Jacobi fiber surfaces that bound it or approximated with the number of regular points in that volume (similarly to the scalar case~\cite{Schneider2008}).

\section{Maybe this should go in method}

Our pipeline ranks sheets at each timestep by range-space area and retains the top $N$ sheets for tracking. Sheets with no associated regular domain vertices are retained, rather than discarded, because they may still represent meaningful range-space structures. For vertex-overlap measures, such sheets are assigned zero overlap.

\myparagraph{Sheet tracking problem.} Given a time varying sequence of bivariate fields
\[
F_0,F_1,\ldots,F_T,
\]
let $\Sheets_i$ denote the set of sheets extracted from the Reeb space of $F_i$. The goal is to identify how sheets evolve between consecutive timesteps by estimating candidate correspondences between sheets in $\Sheets_i$ and $\Sheets_{i+1}$. A correspondence indicates that two sheets may represent the same evolving bivariate feature, or related parts of a feature before or after a topological change.

We do not enforce a one-to-one matching. Reeb space sheets may persist, split, merge, appear, disappear, or change substantially in shape and support. Therefore, for each adjacent pair of timesteps, we construct a weighted bipartite graph whose nodes are sheets from $\Sheets_i$ and $\Sheets_{i+1}$, and whose edge weights encode sheet similarity according to one or more correspondence measures.

\section{Method}
This section describes our pipeline for tracking Reeb space sheets in time-varying bivariate data. Starting from a sequence of volumetric datasets with two scalar fields, we compute the Reeb space at each timestep and extract the most prominent sheets. We then generate visual representations of their range-space geometry and corresponding spatial fiber-surface context. To track sheets over time, we compare sheets from adjacent timesteps using complementary similarity measures in the range and domain. The resulting correspondences define a temporal graph, which we visualize and analyze to identify persistent sheets, weak continuations, and candidate transition intervals.

\subsection{Input and Preprocessing}

\begin{figure*}[t]
  \wideplaceholder{1.55in}{Figure placeholder: pipeline overview. Show input
  time-varying bivariate fields, Reeb space computation, top-sheet extraction,
  sheet and fiber-surface rendering, pairwise sheet matching, event analysis,
  and Sankey-style visual exploration.}
  \caption{Overview of the proposed Reeb space sheet tracking pipeline.}
  \label{fig:pipeline}
\end{figure*}

The input to our method is a time-ordered sequence of volumetric datasets. Each timestep contains two scalar fields defined on a common spatial domain, forming a bivariate field $F_i=(f_i,g_i)$. We assume that the field pair is consistently defined across the sequence so that the corresponding scalar quantities are comparable over time.
\Mohit{Mention how is sheet area computed?}
For each timestep, we compute the Reeb space of the bivariate field and extract its sheet representation. For each sheet, this representation stores its geometric footprint in the range, its range-space area, and the associated set of regular domain vertices, when available. Some sheets may have no associated regular vertices due to input sampling, or the structure of the computed sheet representation. We retain such sheets because they may still represent meaningful range-space geometry. The extracted sheet information provides the basis for sheet ranking, visualization, and correspondence estimation.




Only timesteps for which a valid Reeb space sheet representation is obtained are included in the tracking stage. Timesteps where the Reeb space computation fails, for example due to high degeneracy, are excluded from the temporal correspondence graph.



\subsection{Top-Sheet Selection}

Let $\Sheets_i^\mathrm{all}$ denote the set of all sheets extracted at timestep $i$. We rank these sheets by their range-space area and retain the $N$ largest sheets:
\[
\Sheets_i = \operatorname{TopN}(\Sheets_i^\mathrm{all}, N).
\]
This selection focuses the analysis on dominant bivariate features while keeping the temporal graph manageable. The value of $N$ is a user-defined parameter and can be varied to study the sensitivity of the analysis to smaller or less prominent sheets.


\begin{figure}[t]
  \placeholder{1.35in}{Figure placeholder: one timestep with the full
  Reeb space sheet collection, top sheets highlighted, and a zoomed example of
  one sheet footprint in range space.}
  \caption{A Reeb space sheet is represented by its footprint in the bivariate
  range and by its associated domain support.}
  \label{fig:sheet-definition}
\end{figure}

\subsection{Sheet and Fiber-Surface Renderings}
For interactive analysis, we precompute visual representations of the selected sheets and their corresponding spatial context. Each selected sheet is represented by its footprint as a polygonal region in the bivariate range. These range-space footprints are used for both visualization and shape-based comparison across timesteps. Each sheet is rendered with a uniform color so that the image emphasizes its geometry. The full Reeb space is shown as faint context, while the selected sheet remains visually dominant.

To support meaningful visual comparison over time, all sheet renderings use common range-space bounds and a fixed image size. This ensures that apparent changes in sheet size and position reflect changes in the data, rather than independent rescaling of each rendered image.

We also precompute fiber-surface renderings for selected sheet and timestep combinations. These renderings provide spatial context by showing how a selected range-space sheet is embedded in the domain. The specific fiber surfaces are chosen according to the dataset and analysis task, and are stored for inspection alongside the corresponding Reeb space sheets. Together, the sheet and fiber-surface renderings help relate temporal sheet correspondences to the underlying volumetric structures.

\subsection{Sheet Correspondence Measures}
For each adjacent timestep pair, we estimate candidate correspondences between sheets using similarity measures defined in the domain, in the range, or as a combination of both. Domain-space similarity captures whether two sheets are associated with overlapping spatial regions. Range-space similarity captures whether their geometric footprints in the bivariate range are similar. The hybrid score combines both cues for exploratory analysis.

\subsubsection{Domain-Space Sheet Similarity}

Domain-space similarity compares the spatial regions associated with sheets across adjacent timesteps. Let $V_\Sheet$ and $V_{\Sheet'}$ denote the sets of regular domain vertices associated with a source sheet $\Sheet$ and a target sheet $\Sheet'$, respectively. We first compute their raw spatial overlap as
\[
O(\Sheet,\Sheet') = |V_\Sheet \cap V_{\Sheet'}|.
\]
A positive overlap indicates that the two sheets share part of their spatial region and are therefore candidates for temporal correspondence. To account for changes in sheet size, we compute two directional normalized overlap scores:
\[
P_\mathrm{source}(\Sheet,\Sheet') =
\begin{cases}
\dfrac{|V_\Sheet \cap V_{\Sheet'}|}{|V_\Sheet|}, & |V_\Sheet| > 0,\\
0, & |V_\Sheet| = 0,
\end{cases}
\]
and
\[
P_\mathrm{target}(\Sheet,\Sheet') =
\begin{cases}
\dfrac{|V_\Sheet \cap V_{\Sheet'}|}{|V_{\Sheet'}|}, & |V_{\Sheet'}| > 0,\\
0, & |V_{\Sheet'}| = 0.
\end{cases}
\]
The source-normalized score $P_\mathrm{source}$ measures how much of the source sheet's spatial region overlaps the target sheet, while $P_\mathrm{target}$ measures how much of the target sheet's spatial region overlaps the source sheet. These two scores are directional. Together, they help distinguish temporal behaviors such as continuation, growth, shrinkage, merging, splitting, and partial overlap. We summarize the directional overlaps using
\[
P_\mathrm{max}(\Sheet,\Sheet') =
\max(P_\mathrm{source}(\Sheet,\Sheet'),
P_\mathrm{target}(\Sheet,\Sheet')).
\]
This score treats containment as strong evidence of domain continuity. For example, if a small source sheet is fully contained in a larger target sheet, then $P_\mathrm{source}$ is high even though the target sheet contains additional vertices. This behavior is useful for identifying continuations involving growth, shrinkage, merging, or absorption. A possible alternative is the symmetric Jaccard overlap,
\[
P_\mathrm{Jaccard}(\Sheet,\Sheet') =
\frac{|V_\Sheet \cap V_{\Sheet'}|}
{|V_\Sheet \cup V_{\Sheet'}|}.
\]
However, Jaccard overlap penalizes size imbalance between sheets. This is undesirable in our setting because containment relationships can indicate meaningful temporal behavior. The Jaccard score would be low when the larger target sheet contains many vertices outside the smaller source sheet, even if the source sheet is fully contained within it. We therefore use $P_\mathrm{max}$ as the domain-space similarity score, while retaining $P_\mathrm{source}$ and $P_\mathrm{target}$ to interpret the direction of overlap.


\subsubsection{Range-Space Sheet Similarity}

For each adjacent timestep pair $(i,i+1)$, we compare every source sheet $\Sheet\in\Sheets_i$ with every target sheet $\Sheet'\in\Sheets_{i+1}$. Range-space similarity measures how similar two sheets are as geometric objects in the bivariate range. These measures complement the domain-space overlap described above: two sheets may occupy similar regions in the range even if their associated spatial regions differ, and conversely, sheets with overlapping domain regions may change their range-space shape.

The primary range-space measure is shape intersection-over-union. Let $M_\Sheet$ and $M_{\Sheet'}$ denote binary masks representing the range-space footprints of two sheets under a common range-space coordinate system. We define
\[
\operatorname{shapeIoU}(\Sheet,\Sheet') =
\frac{|M_\Sheet\cap M_{\Sheet'}|}{|M_\Sheet\cup M_{\Sheet'}|}.
\]
This score measures direct overlap between sheet shapes in the range. It is the most informative range-space measure when the goal is to identify sheets that preserve similar geometric footprints across time.

We also compute a bounding-box overlap as a coarser geometric cue. Let $B_\Sheet$ and $B_{\Sheet'}$ be the axis-aligned bounding boxes of the two sheets in range space. The bounding-box IoU is
\[
\operatorname{bboxIoU}(\Sheet,\Sheet') =
\frac{\operatorname{area}(B_\Sheet\cap B_{\Sheet'})}
{\operatorname{area}(B_\Sheet\cup B_{\Sheet'})}.
\]
This measure captures coarse positional and scale agreement. It is less precise than shape IoU because it ignores the detailed sheet geometry, but it can provide useful supporting evidence when sheets have similar extent in the range.

The area-ratio score compares the sizes of two sheets:
\[
\operatorname{areaRatio}(\Sheet,\Sheet') =
\frac{\min(A_\Sheet,A_{\Sheet'})}{\max(A_\Sheet,A_{\Sheet'})},
\]
where $A_\Sheet$ and $A_{\Sheet'}$ are their range-space areas. This score is high when the two sheets have comparable area. It does not encode where the sheets are located in the range and is therefore used only as a size-consistency cue.

Finally, we use centroid similarity as a weak positional measure. Let $c_\Sheet$ and $c_{\Sheet'}$ be the range-space centroids of the two sheets, and let $D$ be the diagonal length of the common range-space bounds. We define
\[
\operatorname{centroidSim}(\Sheet,\Sheet') =
\exp\left(-\frac{\|c_\Sheet-c_{\Sheet'}\|}{D}\right)
\]
This score decreases smoothly as the centroid distance increases. It provides a simple measure of range-space proximity, but it is not sufficiently discriminative on its own because sheets with different shapes may have nearby centroids.

When a single range-space correspondence score is required, we combine these measures using a weighted sum:
\[
\begin{aligned}
\Score(\Sheet,\Sheet') ={}&
w_1\,\operatorname{shapeIoU}(\Sheet,\Sheet') \\
&+w_2\,\operatorname{areaRatio}(\Sheet,\Sheet') \\
&+w_3\,\operatorname{bboxIoU}(\Sheet,\Sheet') \\
&+w_4\,\operatorname{centroidSim}(\Sheet,\Sheet'),
\end{aligned}
\]
where $w_1,w_2,w_3,w_4\geq 0$ are user-defined weights satisfying
\[
w_1+w_2+w_3+w_4=1.
\]
The weights control the relative influence of shape overlap, area consistency, coarse extent overlap, and centroid proximity. In our analysis, shape IoU is treated as the main range-space correspondence cue, while the remaining terms provide supporting geometric evidence. We also inspect the individual scores alongside domain-space similarity to distinguish stable continuations from cases where a sheet preserves its range-space geometry but changes its spatial region, or preserves its spatial region while changing its range-space footprint.

\subsubsection{Hybrid Domain--Range Similarity}

  The pipeline also supports a hybrid similarity score that combines domain-space
  overlap with range-space sheet similarity. This score is intended for
  interactive exploration when the analyst wants to emphasize both spatial support
  continuity in the input domain and geometric continuity in the bivariate range.

  Let $P_\mathrm{domain}(\Sheet,\Sheet')$ denote the selected domain-overlap
  metric. In the current viewer, the default choice is

  \[
  P_\mathrm{domain}(\Sheet,\Sheet') =
  P_\mathrm{max}(\Sheet,\Sheet').
  \]

  Let $C_\mathrm{range}(\Sheet,\Sheet')$ denote the selected range-space
  combined score, computed from range-space metrics such as shape IoU, area
  ratio, bounding-box IoU, and centroid similarity. Since the domain-overlap
  metrics can have different scales, the selected domain metric is normalized by
  its maximum value over the current dataset

  \[
  \widehat{P}_\mathrm{domain}(\Sheet,\Sheet') =
  \frac{P_\mathrm{domain}(\Sheet,\Sheet')}
  {\max_{\Sheet_a,\Sheet_b} P_\mathrm{domain}(\Sheet_a,\Sheet_b)}.
  \]

  The hybrid score is then

  \[
  C_\mathrm{hybrid}(\Sheet,\Sheet') =
  \frac{
  w_\mathrm{domain}\widehat{P}_\mathrm{domain}(\Sheet,\Sheet')
  +
  w_\mathrm{range}C_\mathrm{range}(\Sheet,\Sheet')
  }{
  w_\mathrm{domain}+w_\mathrm{range}
  },
  \]

  where $w_\mathrm{domain}\geq 0$ and $w_\mathrm{range}\geq 0$ are
  user-adjustable weights.

  This hybrid score is not used as the primary metric in the current analysis
  tables because we wanted to keep domain-space and range-space behavior
  separable. However, it is available in the viewer as an exploratory option. It
  can be useful when a user wants to suppress matches that are strong only in
  range space but have weak domain support, or conversely to highlight matches
  that are supported by both similar sheet geometry and shared domain vertices.

\subsection{Temporal Graph and Sankey Layout}
After computing sheet-to-sheet similarities between adjacent timesteps, we construct a temporal correspondence graph for visualization. Each selected sheet is represented as a node, and each candidate correspondence between sheets in consecutive timesteps is represented as an edge. For a timestep pair $(i,i+1)$, an edge connects $\Sheet\in\Sheets_i$ to $\Sheet'\in\Sheets_{i+1}$ when the selected similarity measure assigns a valid correspondence score to the pair. In domain-space mode, edges are created only for sheet pairs with nonzero spatial overlap. In range-space mode, edges are defined for all compared sheet pairs using their range-space similarity. In hybrid mode, edge weights combine the selected domain- and range-space scores.

We visualize this temporal graph using a Sankey-based sheet evolution summary. Timesteps are arranged from left to right, and the selected sheets at each timestep are shown as vertical nodes. Links between adjacent timestep columns represent candidate correspondences. Link thickness and opacity encode the selected similarity score, allowing strong continuations, weak matches, and possible split or merge events to be inspected in a single temporal view.


\myparagraph{Node height.} In the current view, node height encodes the number of regular domain vertices associated with each sheet. Sheets with more associated vertices are therefore drawn taller than sheets with fewer vertices. Sheets with no associated regular vertices are still retained and assigned a small minimum height so that they remain visible and selectable.

Node heights are normalized within the current layout rather than drawn in direct proportion to the raw vertex counts. This normalization is necessary because each timestep may contain many selected sheets, and all nodes must fit within the available vertical space while remaining large enough for interaction. Consequently, node height should be interpreted as a relative visual cue within the Sankey panel, not as an exact quantitative bar chart of vertex counts.

Small differences in node height may occur across overlap, range-shape, and hybrid views because each view may display a different subset of links after metric selection and thresholding. The layout prioritizes readability of the visible graph. Therefore, quantitative comparisons of sheet size should rely on the reported vertex counts rather than on apparent node height alone.



\myparagraph{Link width and opacity.} Each visible link encodes the selected correspondence score $m(\Sheet,\Sheet')$, which may be based on domain-space similarity, range-space similarity, or a hybrid of both. To make scores comparable within a view, we normalize each score by the maximum value of the selected metric over the considered sheet pairs:
\[
\widehat{m}(\Sheet,\Sheet') =
\frac{m(\Sheet,\Sheet')}
{\max_{\Sheet_a,\Sheet_b} m(\Sheet_a,\Sheet_b)}.
\]
The normalized value is clamped to the interval $[0,1]$ and used to control the visual prominence of the link. Link width is mapped linearly between a minimum and maximum thickness:
\[
W(\Sheet,\Sheet') =
W_\mathrm{min}
+
\widehat{m}(\Sheet,\Sheet')
\left(W_\mathrm{max}-W_\mathrm{min}\right).
\]
The same normalized score is also used to control link opacity. Thus, stronger correspondences appear thicker and more opaque, while weaker correspondences remain visible but visually less prominent. This encoding allows the viewer to show both dominant continuations and weaker candidate matches in the same temporal graph.


\myparagraph{Threshold Filtering.} Each panel provides an interactive threshold for filtering weak correspondences. The threshold is applied to the normalized selected score $\widehat{m}$. A link is shown only if
\[
\widehat{m}(\Sheet,\Sheet') \geq \tau,
\]
where $\tau\in[0,1]$ is the user-selected threshold. This allows the analyst to suppress weak links and focus on stronger candidate correspondences. Since thresholding is applied after metric normalization, the threshold is meaningful within a selected metric mode, but should not be interpreted as an equivalent physical value across different metrics.

\myparagraph{Range selection.} The interface includes an interactive timestep-range selector above the Sankey panels. This control allows the analyst to restrict the view to one or more temporal windows instead of displaying the full sequence at once. It is particularly useful for long sequences, where showing all timesteps simultaneously can make the Sankey summary too dense to inspect. The selected range determines which timestep columns and adjacent-timestep links are shown. A link is visible only when both its source and target timesteps lie within the selected range. The selector does not modify the underlying similarity scores or recompute the correspondence graph; it only filters the visible temporal window. This allows the analyst to identify interesting intervals in the full overview and then inspect those intervals in greater detail.


\myparagraph{Strongest outgoing link filtering.}
The interface provides an option to show only the strongest outgoing link from each source sheet. When enabled, the selected metric is evaluated for all visible outgoing links from a source node, and only the link with the largest score is retained. Weaker outgoing links from the same source sheet are hidden. This filtering is metric-dependent. In domain-space mode, the retained link is the target sheet with the highest selected domain-overlap score. In range-space mode, it is the target sheet with the highest selected range-space similarity score. Changing the active metric can therefore change which correspondence is shown as the strongest continuation. This option is useful when the candidate graph is dense and the analyst wants to inspect the dominant continuation hypothesis for each sheet. It does not impose a globally optimal or one-to-one matching: multiple source sheets may still select the same target sheet, and hidden links can be restored by disabling the filter.


\myparagraph{Node Ordering.} The view supports multiple ordering modes for arranging sheets within each timestep. The crossing-minimized ordering reduces visual clutter by placing sheets so that strong links produce fewer crossings between adjacent timesteps. Alternative orderings sort sheets by range-space area or by the number of associated domain vertices. These choices affect only the vertical placement of nodes; they do not change the underlying correspondence graph or the computed similarity scores.

\myparagraph{Node Coloring.} The interface supports several node-coloring modes. A solid-color mode uses a single neutral color for all sheets. Area and vertex-count modes color nodes by normalized range-space area or by the number of associated domain vertices, respectively. The interface also supports centroid-based coloring in the bivariate range. In this mode, each sheet is assigned a color according to the location of its range-space centroid using bounds that are fixed across all timesteps. As a result, centroid colors provide a temporally consistent cue for comparing sheet positions in the range. We use two centroid-based color maps. The first interpolates colors from the four corners of the global two-dimensional range. The second uses an axis-based red-blue map, where hue varies with the angular position of the centroid around the origin and saturation increases with distance from the origin. These color maps are intended as qualitative spatial cues for visual analysis. They do not replace the quantitative similarity scores used for correspondence estimation.


\myparagraph{Panels and metrics.} The interface can display multiple Sankey panels side by side. Each panel can be configured with a different correspondence mode, metric, threshold, and vertical scale. This enables direct comparison of domain-space, range-space, and hybrid correspondences. For example, one panel may show links based on domain vertex overlap, while another shows links based on range-space shape IoU or the combined range-space score. Differences between panels help identify cases where sheets preserve similar range-space geometry but have weak spatial overlap, or where they share a spatial region but undergo substantial range-space deformation.


\myparagraph{Interaction and detail views.} The interface connects the temporal overview to detailed visual information. Hovering over or selecting a node opens a detail view containing sheet metadata, the rendered range-space sheet image, and associated fiber-surface images when available. Selecting a link opens a side-by-side comparison of the source and target sheets connected by that correspondence. Clicking an image in the detail view opens an image viewer with linked pan and zoom, allowing corresponding sheet or fiber-surface renderings to be inspected at the same scale. These linked views provide the range-space geometry and spatial context needed to interpret individual matches.

\subsection{Event Diagnostics}

The Sankey-based sheet evolution summary provides an overview of temporal correspondences, but long sequences may still contain many intervals that require closer inspection. We therefore compute an event diagnostic score for each adjacent timestep pair. The score is not intended to be a formal topological event detector. Instead, it highlights intervals where the correspondence structure suggests weak continuation, possible appearance or disappearance of sheets, or ambiguous split and merge behavior.

Let $\Score(\Sheet,\Sheet')$ denote the selected correspondence score between a source sheet $\Sheet\in\Sheets_i$ and a target sheet $\Sheet'\in\Sheets_{i+1}$. For each source sheet, we define its best outgoing score as
\[
B^\mathrm{out}(\Sheet) =
\max_{\Sheet'\in\Sheets_{i+1}} \Score(\Sheet,\Sheet').
\]
This value measures the strength of the best available continuation of $\Sheet$ into the next timestep. Similarly, for each target sheet, we define its best incoming score as
\[
B^\mathrm{in}(\Sheet') =
\max_{\Sheet\in\Sheets_i} \Score(\Sheet,\Sheet').
\]
This value measures the strength of the best available predecessor of (\Sheet') from the previous timestep.

Given a continuity threshold $\theta$, we count source sheets whose best continuation is weak:
\[
W^\mathrm{out}_i(\theta) =
\left|
\left\{
\Sheet\in\Sheets_i
\;:\;
B^\mathrm{out}(\Sheet) < \theta
\right\}
\right|.
\]
We also count target sheets whose best predecessor is weak:
\[
W^\mathrm{in}_i(\theta) =
\left|
\left\{
\Sheet'\in\Sheets_{i+1}
\;:\;
B^\mathrm{in}(\Sheet') < \theta
\right\}
\right|.
\]
Large values of these terms indicate intervals where many sheets lack a strong continuation or predecessor.

We next count possible branching configurations. A source sheet is counted as a possible split if it has at least two above-threshold correspondences in the next timestep:
\[
P^\mathrm{split}_i(\theta) =
\left|
\left\{
\Sheet\in\Sheets_i
\;:\;
\left|
\left\{
\Sheet'\in\Sheets_{i+1}
\;:\;
\Score(\Sheet,\Sheet') \geq \theta
\right\}
\right|
\geq 2
\right\}
\right|.
\]
Similarly, a target sheet is counted as a possible merge if it has at least two above-threshold correspondences from the previous timestep:
\[
P^\mathrm{merge}_i(\theta) =
\left|
\left\{
\Sheet'\in\Sheets_{i+1}
\;:\;
\left|
\left\{
\Sheet\in\Sheets_i
\;:\;
\Score(\Sheet,\Sheet') \geq \theta
\right\}
\right|
\geq 2
\right\}
\right|.
\]
These counts indicate correspondence ambiguity. Such ambiguity may correspond to genuine splitting or merging, but it may also occur when several sheets have similar range-space geometry or domain-space overlap under the selected metric.

Finally, we compute the mean best outgoing score,
\[
\overline{B^\mathrm{out}_i} =
\frac{1}{|\Sheets_i|}
\sum_{\Sheet\in\Sheets_i} B^\mathrm{out}(\Sheet).
\]
This term summarizes the overall continuation strength from timestep $i$ to timestep $i+1$.

The event diagnostic score for the interval $i\rightarrow i+1$ is
\[
\begin{aligned}
E_i(\theta)={}&
W^\mathrm{out}_i(\theta)
+
W^\mathrm{in}_i(\theta)
+
\frac{1}{2}
\left(
P^\mathrm{split}_i(\theta)
+
P^\mathrm{merge}_i(\theta)
\right) \\
&+
\left(1-\overline{B^\mathrm{out}_i}\right)
|\Sheets_i|.
\end{aligned}
\]
The first two terms emphasize sheets with weak continuations or weak predecessors. The split and merge terms emphasize ambiguous branching, but with a smaller weight because such sheets still have above-threshold correspondences. The final term increases the score when the overall best-continuation strength is low.

Large values of $E_i(\theta)$ indicate intervals with weak temporal continuity or ambiguous correspondence structure. We use these intervals as candidates for closer visual inspection in the Sankey summary and linked sheet views. In our analysis, $\theta=0.5$ is used as the default threshold for ranking intervals, and additional thresholds are used to inspect the sensitivity of the diagnostic.



\subsection{Sheet Lifetimes}

In addition to identifying event-like intervals, we compute a simple lifetime diagnostic for individual sheets. The goal is to estimate how long a sheet can be followed through time under a selected correspondence score and continuity threshold.

The lifetime computation uses the same sheet-to-sheet score $\Score(\Sheet,\Sheet')$ used in the event diagnostics. For a sheet $\Sheet_i \in \Sheets_i$, we first identify its best continuation in the next timestep:
\[
\Sheet_{i+1}^* =
\arg\max_{\Sheet'\in\Sheets_{i+1}}
\Score(\Sheet_i,\Sheet').
\]
If the corresponding score satisfies
\[
\Score(\Sheet_i,\Sheet_{i+1}^*) \geq \theta,
\]
the track is extended to $\Sheet_{i+1}^*$. The same greedy step is then repeated from $\Sheet_{i+1}^*$ to the following timestep. The lifetime of the initial sheet is defined as the number of timesteps visited before the best available continuation falls below the threshold or the sequence ends.

This procedure produces a greedy best-continuation track. It does not enforce a global one-to-one assignment, and multiple sheets from earlier timesteps may select the same sheet in a later timestep. We compute lifetimes for the same continuity thresholds used in the event diagnostics. Unless otherwise stated, we report summaries using $\theta=0.5$.

%Mohit: Summarized interaction, if details have to compressed to save space
%\subsection{Visual Encoding and Interaction}

%The unified viewer represents sheets as nodes and candidate correspondences as links. Node color can encode sheet rank, metric value, or range-space centroid position. For centroid-based coloring, we use global range-space extents so that colors remain comparable across timesteps. The viewer supports two centroid color maps: a four-corner interpolation map and an axis-oriented map in which saturation increases with distance from the origin along red--blue directions. Hovering over or selecting a node shows the corresponding sheet image and associated fiber-surface images when available. Hovering over or selecting a link opens a side-by-side comparison of the connected source and target sheets, together with their fiber-surface images. The detail view supports linked pan and zoom, allowing both sides of a comparison to be inspected at the same scale.


\begin{figure*}[t]
  \wideplaceholder{1.75in}{Figure placeholder: unified viewer screenshot.
  Include a Sankey diagram with node/link coloring, a selected link, side-by-side
  sheet images, and side-by-side fiber-surface images with linked zoom.}
  \caption{The visual interface links temporal sheet correspondences with
  range-space sheet images and spatial fiber-surface evidence.}
  \label{fig:viewer}
\end{figure*}


%% if specified like this the section will be committed in review mode
\acknowledgments{}

%\bibliographystyle{abbrv}
\bibliographystyle{abbrv-doi}
%\bibliographystyle{abbrv-doi-narrow}
%\bibliographystyle{abbrv-doi-hyperref}
%\bibliographystyle{abbrv-doi-hyperref-narrow}

\bibliography{bib}
\end{document}
